% =============================================================
% Cost, data, reliability: the minimum vocabulary
% =============================================================

\sectionframe[The minimum vocabulary before Lesson 3]{Cost, data, reliability}

% -------- Subscription vs API, four levers --------
\begin{frame}{Subscription vs API, and what drives the bill}
  \begin{columns}[T,onlytextwidth]
    \column{0.42\textwidth}
      \group{Subscription}
      {\small\color{muted}Flat monthly price. Limits are real but soft and change without notice. Fine for chat and desktop apps.\par}
      \vspace{0.7em}
      \group{API}
      {\small\color{muted}Pay per token. Easy to measure, easy to overspend. Needed when a harness runs many steps on its own.\par}
      \vspace{0.9em}
      {\scriptsize\color{muted}List prices per million tokens, input and output: GPT-6 Luna \$0.10 and \$0.50; GPT-6.1 Sol \$2 and \$10; Claude Opus 5.5 \$4 and \$20; GPT-6 Astra \$10 and \$50.\par}
    \column{0.54\textwidth}
      \group{Four levers of an agent's bill}
      \vspace{0.2em}
      \renewcommand{\arraystretch}{1.4}
      \begin{tabularx}{\linewidth}{>{\bfseries\color{ink}}l >{\small\leavevmode\color{dark}}X}
        Context & every step resends files, history and tool output\\
        Tool calls & more autonomy, more steps\\
        Retries & failed runs, self-checks, sub-agents\\
        Model tier & a 100-fold range between Luna and Astra\\
      \end{tabularx}
      \par\vspace{1.0em}
      {\small\color{ink}\bfseries A chat ignores three of these. An agent pays for all four.\par}
  \end{columns}
  \source{OpenAI and Anthropic API pricing pages, read 26 Oct 2026.}
\end{frame}

% -------- Cost spike --------
\begin{frame}{What a cost spike looks like}
  \vspace{0.5em}
  \begin{center}
  \renewcommand{\arraystretch}{1.3}
  \begin{tabularx}{0.9\linewidth}{>{\bfseries\color{ink}}l X r}
  \toprule
  \color{muted}Workload & \leavevmode\color{muted}Example & \color{muted}Tokens, roughly\\
  \midrule
  Chat question & ``Explain heteroskedasticity in two lines.'' & 0.5k\\
  Long reading & ``Summarise this 30-page paper.'' & 25k\\
  Agent task & ``Rerun all tables and check them.'' & 100k to 200k\\
  Long agent run & ``Replicate Section 4 end to end.'' & 0.5M to 2M\\
  Full context, ten retries & a 1M-token window resent ten times & 10M\\
  \bottomrule
  \end{tabularx}
  \end{center}
  \vspace{0.4em}
  {\small\color{muted}\centering Ten resends of a 1M-token prompt on GPT-6.1 Sol: \$40 at the long-context list price, about \$2 if the context is cached.\par}
  \takeaway{Same window on screen. Four orders of magnitude in cost.}
\end{frame}

% -------- Data risk --------
\begin{frame}{Does it train on my data, and what is risky}
  \begin{columns}[c,onlytextwidth]
    \column{0.52\textwidth}
      \begin{tikzpicture}[y=-1.25cm]
        \foreach \i/\lvl/\ex/\c in {%
          0/High risk/{microdata, administrative or embargoed data, referee reports, IRB-restricted material}/bad,%
          1/Medium/{unpublished drafts, grant ideas, coauthors' notes, interview transcripts}/accent,%
          2/Lower/{public documentation, published papers, teaching material, synthetic data}/faint}{%
          \fill[\c] (0,\i) rectangle ++(0.12,0.95);
          \node[anchor=north west,font=\small\bfseries,text=ink] at (0.3,\i) {\lvl};
          \node[anchor=north west,font=\scriptsize,text=muted,text width=5.6cm] at (0.3,\i+0.38) {\ex};
        }
      \end{tikzpicture}
    \column{0.44\textwidth}
      {\small
      \group{Whether it trains depends on}
      {\color{muted}the plan (consumer or business), the workspace settings, the API retention policy, and every tool the agent calls on the way.\par}
      \vspace{0.8em}
      {\color{ink}\bfseries Settings change. Check them; do not assume them.\par}
      \vspace{0.8em}
      {\color{muted}When in doubt: anonymise, use a synthetic copy, or run locally.\par}}
  \end{columns}
\end{frame}

% -------- Reliability words --------
\begin{frame}{Reliability words you will keep hearing}
  \vspace{0.4em}
  \begin{center}
  \renewcommand{\arraystretch}{1.25}
  \begin{tabularx}{0.9\linewidth}{>{\bfseries\color{ink}}l X}
  \toprule
  Hallucination & a confident output with no basis: a citation, a number, a variable\\
  Context rot & quality falls as the context window fills\\
  Reward hacking & the agent meets the check and misses the goal, for example by editing the test instead of the code\\
  Sycophancy & the model agrees with you when you are wrong\\
  Automation bias & we trust automated output more than it deserves \citep{skitka_1999_automation_bias}\\
  Drift & the same prompt gives different results after a model update\\
  \bottomrule
  \end{tabularx}
  \end{center}
  \takeaway{Lesson 3 is about controlling all six.}
\end{frame}
